%%BLOCK after="What the bifibration question asks"
\begin{dcomment}{development}{p186}{Bifibrations over posets and over bases with cancellation}
For $W=W_\infty$ the essential surjectivity asked for on this page can be proved for a large class of bases, among them all partially ordered sets. The statement is: every locally constant object of $\mathsf{HOT}(S)$, that is, every diagram of spaces on $S$ whose transition maps are weak equivalences, is the diagram of fibres of a functor $X\to S$ that is at once a Grothendieck fibration and a Grothendieck opfibration. The theorem below was not found in the sources consulted; no priority is claimed.

\textbf{Lemma 1.}
\textit{A Grothendieck fibration is an opfibration if and only if each of its transition functors has a left adjoint. The transition functors of a bifibration are therefore weak equivalences, and a bifibration is proper and smooth.}

This is classical: a cocartesian lift of $\alpha:s\to t$ at $x$ is an initial object of the category of lifts of $\alpha$ with source $x$, which is the comma category of $x$ over the transition functor $\alpha^*:X_t\to X_s$. The converse of the last assertion fails. The four-element poset $x<a,b<c$, mapped to $[1]$ by $x\mapsto0$ and $a,b,c\mapsto1$, is proper and smooth over $[1]$ (the lifting posets $\{a,b,c\}$ and $\{x\}$ have a greatest element) and is a fibration (the cartesian lifts end at $x$); it is not an opfibration, since $x$ has no least element above it in the fibre over $1$.

\textbf{Lemma 2 (patches).}
\textit{Let $Y$ be a simplicial set and $\lambda$ an infinite cardinal with $|Y|\leq\lambda$. Let $K_\lambda(Y)$ be the category of the nonempty weakly contractible simplicial sets of cardinality at most $\lambda$ over $Y$, with all maps over $Y$. The functor $j:\Delta/Y\to K_\lambda(Y)$ which regards a simplex of $Y$ as an object over $Y$ is a weak equivalence. Hence $BK_\lambda(Y)$ is naturally weakly equivalent to $|Y|$.}

\textit{Proof.}
For an object $P\to Y$ of $K_\lambda(Y)$ the comma category $j/P$ is $\Delta/P$, whose classifying space has the weak homotopy type of $P$ and is contractible. Quillen's Theorem~A applies; and $B(\Delta/Y)\to|Y|$ is a natural weak equivalence.\hfill$\square$

\textbf{Lemma 3.}
\textit{Let $f:Y\to Z$ be a trivial Kan fibration of simplicial sets of cardinality at most $\lambda$. Then composition with $f$, $f_!:K_\lambda(Y)\to K_\lambda(Z)$, has the right adjoint $f^*$, pullback along $f$.}

\textit{Proof.}
For $Q\to Z$ in $K_\lambda(Z)$ the projection $Q\times_ZY\to Q$ is a trivial Kan fibration, so $Q\times_ZY$ is nonempty and weakly contractible, of cardinality at most $\lambda$. The adjunction is that of the slice categories.\hfill$\square$

\textbf{Theorem.}\ours
\textit{Let $S$ be a small category in which every arrow is a monomorphism, or every arrow is an epimorphism; for instance a partially ordered set. Then every locally constant object of $\mathsf{HOT}(S)$ is represented by the diagram of fibres of a bifibration $X\to S$ with $X$ small.}

\textit{Proof.}
Suppose first that every arrow is a monomorphism, and represent the object by a functor $F:S^{\mathrm{op}}\to\mathbf{sSet}$ whose transition maps are weak equivalences. Let $F\to RF$ be a fibrant replacement for the injective model structure on simplicial presheaves on $S$, in which the weak equivalences are objectwise and the cofibrations are the objectwise monomorphisms. We claim that for every arrow $\alpha:s\to t$ the map $RF(\alpha):RF_t\to RF_s$ is a trivial Kan fibration. Let $h_s\to h_t$ be the map of representable presheaves induced by $\alpha$; it is objectwise injective because $\alpha$ is a monomorphism. A lifting problem of a horn inclusion $\Lambda^n_k\subset\Delta^n$ against $RF(\alpha)$ is, by the Yoneda lemma, a lifting problem of
\[
 (h_t\times\Lambda^n_k)\amalg_{h_s\times\Lambda^n_k}(h_s\times\Delta^n)\longrightarrow h_t\times\Delta^n
\]
against $RF\to*$. At each object this map is a disjoint union of identities of $\Delta^n$ and of horn inclusions; it is an objectwise trivial cofibration, and $RF$ is injectively fibrant. So $RF(\alpha)$ is a Kan fibration; it is a weak equivalence because $F(\alpha)$ is one and $F\to RF$ is an objectwise weak equivalence.

Choose $\lambda$ bounding the values of $RF$. By Lemma~3 the functor $s\mapsto K_\lambda(RF_s)$, with the transition functors $RF(\alpha)_!$, takes values in functors having right adjoints; so $s\mapsto K_\lambda(RF_s)^{\mathrm{op}}$ is a functor on $S^{\mathrm{op}}$ whose transition functors have left adjoints. Its Grothendieck construction $X\to S$ is a fibration, and by Lemma~1 a bifibration. By Lemma~2, and because a category and its opposite have the same classifying space naturally, its diagram of fibres is naturally weakly equivalent to $RF$, hence to $F$.

If every arrow is an epimorphism, the covariant representables give objectwise injections $S(t,-)\to S(s,-)$, and the same argument applies to a covariant model of the local system, with Lemma~3 used on the patch categories themselves.\hfill$\square$

Two further cases hold over every small base $S$, by simpler arguments which are not reproduced here: when the values of the local system are 1-types (take fundamental groupoids objectwise; equivalences of groupoids have adjoints), and when $BS$ is a 1-type (pull back from the groupoid obtained by inverting all arrows of $S$; this contains the case of a groupoid treated above, and the bases with an initial or a terminal object).

\textbf{What remains.}
For a general base the question is undecided. The first test case is a monoid $M$ with five elements, the identity and a rectangular band $\{e_{ij}\}_{i,j\in\{1,2\}}$ with $e_{ij}e_{kl}=e_{il}$, whose classifying space has the homotopy type of $S^2$, with the local system of the path fibration (fibre $\Omega S^2$). The method above cannot reach it. A strict action of a monoid on a simplicial set in which an idempotent $e$ acts by a trivial Kan fibration has $e$ acting as the identity, because a surjective idempotent map is the identity; so a strict diagram with trivial fibrations as transitions on a monoid generated by idempotents is constant, and a bifibration for the test case must come from a pseudofunctor that is not strict. The second half of the question of this page, whether the comparison of localizations is an equivalence, is not addressed by the theorem.
\end{dcomment}

%%APPEND title="Why alternating completion works at the limit"
For $W_\infty$ the two assertions hold without the axiom W(6 bis) as a separate hypothesis: the lifting categories of a filtered colimit are the filtered colimits of the lifting categories, and filtered colimits of weakly contractible simplicial sets are weakly contractible. The limit is proper and smooth; it need not be a bifibration, even when every stage is one. Let $X_n$ be the chain $0<1<\dots<n<t$ over $[1]$, with $t$ over $1$ and the rest over $0$. Each $X_n\to[1]$ is a bifibration (the cartesian lift of $0\to1$ ends at $n$). The colimit of the inclusions $X_n\subset X_{n+1}$ has the fibre $\mathbf N$ over $0$, and a cartesian lift with target $t$ would be a greatest natural number. This is why the alternating construction $\Phi_\infty$ yields perfect objects and not bifibrations.
